# 5. Conclusion and Future Work

  

尽管现有覆盖路径规划方法已经能够处理带孔洞、非凸边界和多障碍的二维空间，SatAOI 也能够界定摆臂抹平机器人的可作业区域，但在复杂施工地图中，仍缺少一种把可作业区域进一步转化为完整、连续且满足机器人几何与运动约束的车体—摆臂覆盖轨迹的方法。针对这一缺口，本文提出了一种面向已知二维障碍环境的分层覆盖路径规划方法：将复杂自由空间分解为适合摆扫的矩形作业单元，通过候选选择、补洞、连通修复和全局访问组织连续作业，并在轨迹生成与评价中显式保留车体、摆臂和圆形抹盘的几何关系。与把工具等效为固定宽度的传统模型不同，本文根据抹盘沿连续运动形成的扫掠区域评价覆盖，并将车体线速度与摆臂角速度的匹配纳入这一过程，从而区分静态几何可达与动态执行覆盖。

  

三组实验从覆盖能力、地图适应性和速度匹配三个方面验证了该规划链条。在 `map_test1` 和障碍更密集的 `parkingGraph` 中，本文方法分别达到 94.11% 和 95.73% 的覆盖率，并通过完整路径检查；代表固定宽度 CPP 综合 SOTA 水平的 Fields2Cover 虽生成了标称时间更短的路线，但连接段分别产生 4.51 m² 和 20.00 m² 的越界，因而不能视为可执行方案。候选扫描实验表明，row-then-col 在五张地图上取得最高的平均全轨迹覆盖率和最低的平均总时长，适合作为对地图方向变化更稳健的默认策略，但单方向策略在部分规则场景或单项指标上仍可能更优。速度耦合实验进一步表明，车速提高或摆臂角速度相对不足会增大相邻摆扫间距；即使几何路径和静态可达范围不变，覆盖率仍会下降。由此，本文工作的核心不是以矩形分解替换已有 CPP 的某个环节，而是将复杂地图分解、车体—摆臂协同轨迹生成和连续扫掠评价连接为同一规划流程。

  

本文仍处于从几何规划走向真实施工应用的中间阶段。当前方法采用启发式矩形选择和访问排序，不保证全局最优；输入被限定为已知静态二维地图，速度模型也尚未包含完整底盘动力学、抹盘压力、表面质量反馈、定位误差和动态障碍。后续研究将据此开展分区与访问次序的联合优化，引入在线感知和局部重规划，把动力学与施工质量纳入轨迹评价，并在真实摆臂抹平机器人上验证规划、控制和表面处理效果。在实现完全自主作业之前，本文方法也可作为施工地图分区、覆盖路线生成和车体—摆臂速度配置的决策基础，辅助操作者减少依赖经验的路径选择，并为摆臂抹平机器人进一步提高自主化程度提供完整的规划起点。


